%==================================================================
% SAMPUL MUKA / SOFTCOVER
% Acuan: Pedoman Lampiran 4
%
% Ingin sampul berwarna sesuai program studi? Hapus tanda % pada
% empat baris tikz di bawah (butuh dua kali kompilasi).
%==================================================================
\begin{titlepage}
\thispagestyle{empty}%
\centering

%--- Latar berwarna (opsional) -------------------------------------
% \definecolor{warnasampul}{HTML}{\warnasampul}
% \begin{tikzpicture}[remember picture, overlay]
%   \fill[warnasampul] (current page.south west) rectangle (current page.north east);
% \end{tikzpicture}%
%-------------------------------------------------------------------

\vspace*{0.5cm}

%--- Judul: 14 pt, tebal, kapital, spasi 1 -------------------------
\begin{spacing}{1}
\begin{center}
{\fontsize{14}{16.8}\selectfont\bfseries \MakeUppercase{\judulid}\par}
\end{center}
\end{spacing}

\vspace{4\baselineskip}

%--- Jenis dokumen: 12 pt, tebal, spasi 1 --------------------------
\begin{spacing}{1}
\begin{center}
{\fontsize{12}{14.4}\selectfont\bfseries \juduljenis\par}
\end{center}
\end{spacing}

\vspace{2\baselineskip}

%--- Lambang Poliwangi 4 cm x 4 cm --------------------------------
\includegraphics[width=4cm,height=4cm,keepaspectratio]{logo-poliwangi.png}

\vspace{3\baselineskip}

%--- Nama dan NIM: 12 pt, tebal, spasi 1 --------------------------
\begin{spacing}{1}
\begin{center}
{\fontsize{12}{14.4}\selectfont\bfseries
 Oleh:\par
 \penulis\par
 NIM. \nim\par}
\end{center}
\end{spacing}

\vspace{2\baselineskip}

%--- Program studi dan tahun: 16 pt, tebal, spasi 1 ----------------
\begin{spacing}{1}
\begin{center}
{\fontsize{16}{19.2}\selectfont\bfseries
 PROGRAM STUDI \MakeUppercase{\prodi}\par
 \MakeUppercase{\institusi}\par
 \tahun\par}
\end{center}
\end{spacing}

\end{titlepage}
